\documentclass{beamer}
\usepackage{../tut-slides}
\usepackage{../mathoperators}
\usepackage{../fs}

\usepackage{csquotes}

\usepackage{amsmath,amssymb}
%\usepackage{enumerate}
\usepackage[normalem]{ulem}
\newcommand{\labelitemi}{\raisebox{1pt}{\scalebox{.9}{$\blacktriangleright$}}}
\newcommand{\labelitemii}{$\vartriangleright$}
\newcommand{\labelitemiii}{--}

\usepackage{booktabs}
\usepackage{tabularx}

\newcommand{\tuple}[1]{\langle{#1}\rangle}
\newcommand{\simquot}[1]{#1/_{\!\!{\sim}}}
\newcommand{\bigmid}{\,\Big|\,}

%%%%%%%%%%%%%%%%%%%%%%%%%%%%%%%%%%%%%%%%%%%%%%%%%%%%%%%%%%%%%%%%%%%%%%%


\begin{document}	
	\title{Formale Systeme}
	\subtitle{Übung 4}
	\author{Eric Kunze}
	\email{eric.kunze@tu-dresden.de}
	\city{TU Dresden}
	\date{\today}
%	\institute{Lehrstuhl für Grundlagen der Programmierung}
	\titlegraphic{\includegraphics[width=2cm]{../TUD-white.pdf}}

	\maketitle

	\section{Aufgabe 1 \& 2: \\ \itshape Reguläre Ausdrücke}
	
	\begin{frame} \frametitle{Reguläre Ausdrücke}
		\small		
		\defbox{Die Menge der \textbf{regulärer Ausdrücke} über einem Alphabet $\Sigma$ ist induktiv wie folgt definiert:
			\begin{itemize}
				\item $\emptyset$ ist ein regulärer Ausdruck
				\item $\epsilon$ ist ein regulärer Ausdruck
				\item $\Sterm{a}$ ist ein regulärer Ausdruck für jedes $\Sterm{a}\in\Sigma$
				\item Wenn $\alpha$ und $\beta$ reguläre Ausdrücke sind,\\dann sind auch $(\alpha\beta)$, $(\alpha\mid\beta)$ und $(\alpha)^\ast$ reguläre Ausdrücke
			\end{itemize}
		}
	
		\defbox{Die \textbf{Sprache eines regulären Ausdrucks} $\alpha$ wird mit $\Slang{L}(\alpha)$
			bezeichnet und rekursiv definiert:
			
			\begin{minipage}{\dimexpr0.5\linewidth-\fboxrule-\fboxsep}
				\begin{itemize}
					\item $\Slang{L}(\emptyset)=\emptyset$
					\item $\Slang{L}(\epsilon)=\{\epsilon\}$
					\item $\Slang{L}(\Sterm{a})=\{\Sterm{a}\} \quad \forall \Sterm{a}\in\Sigma$
				\end{itemize}
			\end{minipage}
			\begin{minipage}{\dimexpr0.5\linewidth-\fboxrule-\fboxsep}
				\begin{itemize}
					\item $\Slang{L}((\alpha\beta))=\Slang{L}(\alpha)\circ\Slang{L}(\beta)$
					\item $\Slang{L}((\alpha\mid\beta))=\Slang{L}(\alpha)\cup\Slang{L}(\beta)$
					\item $\Slang{L}((\alpha)^\ast)=\Slang{L}(\alpha)^\ast$
				\end{itemize}
			\end{minipage}
		}
	\end{frame}


	\begin{frame} \frametitle{Aufgabe 1}
		\small
		Gegeben sind das Alphabet $\Sigma=\{a,b,c\}$ und die Sprache
		\begin{equation*}
			L= \{w\in \Sigma^\ast \mid
			\begin{array}[t]{@{\,}l}
				\mbox{es gibt } u,v\in \Sigma^\ast \mbox{
					mit } w = u babc v \mbox{ und }\\
				\mbox{es gibt } u,v\in \Sigma^\ast \mbox{
					mit } w = u ccc v \mbox{ und }\\
				\mbox{es gibt kein } u \in \Sigma^\ast
				\mbox{ mit } w = au\} .
			\end{array}
		\end{equation*}	
		Geben Sie für $L$ einen regulären Ausdruck $r$ mit $L=L(r)$ an.
	\end{frame}

	\begin{frame}
		\small
		\textbf{Idee}: 
		\begin{itemize}
			\item die ersten beiden Bedingungen beschreiben zusammen die \enquote{Wortmitte}
			\item die Teile $ccc$ und $babc$ können in drei verschiedenen Kombinationen auftreten
			\item die dritte Bedingung schließt ein $a$ am Anfang aus
		\end{itemize} 
		\pause
	
		\textbf{Lösung}:
		\begin{align*}
			r \, = \enskip
				&\Big((b \mid c) (a \mid b \mid c)^\ast \Big)^\ast \\
				&\Big( (babc \, (a \mid b \mid c)^\ast \, ccc) 
				\bigmid ccc \, (a \mid b \mid c)^\ast \, babc 
				\bigmid babccc \Big) \\
				&\Big(a \mid b \mid c \Big)^\ast
		\end{align*}
	\end{frame}
	
	\begin{frame} \frametitle{Aufgabe 2}
		\small
		Beweisen Sie die folgenden Gleichungen für reguläre Ausdrücke $r$,
		$s$ und $t$.
		
		\textbf{Erinnerung}: $r\equiv s$ bedeutet $L(r) = L(s)$
		
		\begin{enumerate}[a)]
			\item $r\mid s\equiv s\mid r$
			\item $(r\mid s)\mid t\equiv r\mid (s\mid t)$
			\item $(rs)t\equiv r(st)$
			\item $r(s\mid t)\equiv rs\mid rt$
			\item $\emptyset^\ast\equiv \varepsilon$
			\item $(r^\ast)^\ast\equiv r^\ast$
			\item $r^\ast\equiv rr^\ast\mid \varepsilon$
			\item $(\varepsilon \mid r)^\ast\equiv r^\ast$
		\end{enumerate}
	\end{frame}

	\section{Aufgabe 3: \\ \itshape RegExp $\to$ NFA}

	\begin{frame}\frametitle{RegExp $\to$ NFA: kompositionelle Methode}
		\small
		Für einen Ausdruck $\alpha$ definieren wir rekursiv den $\epsilon$-NFA $\Smach{M}(\alpha)$:
		
		\textbf{Grundfälle:}
		\begin{itemize}
			\item Wenn $\alpha=\emptyset$ dann $\Smach{M}(\alpha)={}$ \scalebox{1}{%
				\begin{tikzpicture}[baseline={(q1.base)}]
					% \draw[help lines] (0,0) grid (7,2);
					\node (q1) [circle,draw=black,thick] at (0,0) {$A$};
					%
					\path[->,line width=0.5mm](-0.7,0) edge (q1);
			\end{tikzpicture}}
			%
			\item Wenn $\alpha=\epsilon$ dann $\Smach{M}(\alpha)={}$ \scalebox{1}{%
				\begin{tikzpicture}[baseline={(q1.base)}]
					% \draw[help lines] (0,0) grid (7,2);
					\node (q1) [circle,draw=black,thick, double] at (0,0) {$A$};
					%
					\path[->,line width=0.5mm](-0.7,0) edge (q1);
			\end{tikzpicture}}
			\item Wenn $\alpha=\Sterm{a}$ dann $\Smach{M}(\alpha)={}$ \scalebox{1}{%
				\begin{tikzpicture}[baseline={(q1.base)}]
					% \draw[help lines] (0,0) grid (7,2);
					\node (q1) [circle,draw=black,thick] at (0,0) {$A$};
					\node (q2) [circle,draw=black,thick,double] at (2,0) {$B$};
					%
					\path[->,line width=0.5mm](-0.7,0) edge (q1);
					\path[->,line width=0.5mm](q1) edge node[above] {\Sterm{a}} (q2);
			\end{tikzpicture}}
		\end{itemize}
		
		\textbf{Rekursive Fälle:}
		Wir bezeichnen mit \bluealert{$\textsf{elim}_\epsilon(\Smach{M})$} den NFA, der aus einem $\epsilon$-NFA $\Smach{M}$ durch Eliminierung der $\epsilon$-Übergänge entsteht.
		\begin{itemize}
			\item Wenn $\alpha=(\beta\gamma)$ dann $\Smach{M}(\alpha)=\textsf{elim}_\epsilon(\Smach{M}(\beta)\odot\Smach{M}(\gamma))$
			\item Wenn $\alpha=(\beta\mid\gamma)$ dann $\Smach{M}(\alpha)=\Smach{M}(\beta)\oplus\Smach{M}(\gamma)$
			\item Wenn $\alpha=(\beta)^\ast$ dann $\Smach{M}(\alpha)=\textsf{elim}_\epsilon(\Smach{M}(\beta)^\ast)$
		\end{itemize}
	\end{frame}
	
	\begin{frame}\frametitle{RegExp $\to$ NFA: Explizite Konstruktion}
		\small
		\textbf{Gegeben}: regulärer Ausdruck $\alpha$ ohne $\emptyset$
		
		\textbf{Initialisierung}: $\Smach{M}_{\alpha}={}$ \scalebox{1}{%
			\begin{tikzpicture}[baseline={(q1.base)}]
				% \draw[help lines] (0,0) grid (7,2);
				\node (q1) [circle,draw=black,thick] at (0,0) {$A$};
				\node (q2) [circle,draw=black,thick,double] at (2,0) {$B$};
				%
				\path[->,line width=0.5mm](-0.7,0) edge (q1);
				\path[->,line width=0.5mm](q1) edge node[above] {$\alpha$} (q2);
		\end{tikzpicture}}
		
		Solange es in $\Smach{M}_{\alpha}$ Übergänge $q\stackrel{\beta}{\to}p$ gibt, die mit einem Ausdruck $\beta\notin \{\epsilon\}\cup\Sigma$ beschriftet sind, wende eine der 
		folgenden Regeln an:
		%
		\begin{itemize}
			\item  Ersetze \begin{tikzpicture}[baseline={(q1.base)}]
				% \draw[help lines] (0,0) grid (7,2);
				\node (q1) [circle,draw=black,thick] at (0,0) {$q$};
				\node (q2) [circle,draw=black,thick] at (1.5,0) {$p$};
				%
				% \path[->,line width=0.5mm](-0.7,0) edge (q1);
				\path[->,line width=0.5mm](q1) edge node[above] {$\enskip (\gamma_1\gamma_2) \enskip $} (q2);
			\end{tikzpicture}
			durch 
			\begin{tikzpicture}[baseline={(q1.base)}]
				% \draw[help lines] (0,0) grid (7,2);
				\node (q1) [circle,draw=black,thick] at (0,0) {$q$};
				\node (r) [circle,draw=black,thick] at (1.5,0) {$r$};
				\node (q2) [circle,draw=black,thick] at (3,0) {$p$};
				%
				% \path[->,line width=0.5mm](-0.7,0) edge (q1);
				\path[->,line width=0.5mm](q1) edge node[above] {$\gamma_1$} (r);
				\path[->,line width=0.5mm](r) edge node[above] {$\gamma_2$} (q2);
			\end{tikzpicture}
			%
			\item Ersetze \begin{tikzpicture}[baseline={(q1.base)}]
				% \draw[help lines] (0,0) grid (7,2);
				\node (q1) [circle,draw=black,thick] at (0,0) {$q$};
				\node (q2) [circle,draw=black,thick] at (2,0) {$p$};
				%
				% \path[->,line width=0.5mm](-0.7,0) edge (q1);
				\path[->,line width=0.5mm](q1) edge node[above,xshift=-3pt] {$(\gamma_1\mid\gamma_2)$} (q2);
			\end{tikzpicture}
			durch 
			\begin{tikzpicture}[baseline={(q1.base)}]
				% \draw[help lines] (0,0) grid (7,2);
				\node (q1) [circle,draw=black,thick] at (0,0) {$q$};
				\node (q2) [circle,draw=black,thick] at (2,0) {$p$};
				%
				% \path[->,line width=0.5mm](-0.7,0) edge (q1);
				\path[->,line width=0.5mm,bend left=20](q1) edge node[above] {$\gamma_1$} (q2);
				\path[->,line width=0.5mm, bend right=20](q1) edge node[above] {$\gamma_2$} (q2);
			\end{tikzpicture}
			%
			\item  Ersetze \begin{tikzpicture}[baseline={(q1.base)}]
				% \draw[help lines] (0,0) grid (7,2);
				\node (q1) [circle,draw=black,thick] at (0,0) {$q$};
				\node (q2) [circle,draw=black,thick] at (1.5,0) {$p$};
				%
				% \path[->,line width=0.5mm](-0.7,0) edge (q1);
				\path[->,line width=0.5mm](q1) edge node[above] {$(\gamma)^\ast$} (q2);
			\end{tikzpicture}
			durch 
			\begin{tikzpicture}[baseline={(q1.base)}]
				% \draw[help lines] (0,0) grid (7,2);
				\node (q1) [circle,draw=black,thick] at (0,0) {$q$};
				\node (r) [circle,draw=black,thick] at (1.5,0) {$r$};
				\node (q2) [circle,draw=black,thick] at (3,0) {$p$};
				%
				\path[->,line width=0.5mm](-0.7,0) edge (q1);
				\path[->,line width=0.5mm](q1) edge node[above] {$\epsilon$} (r);
				\path[->,line width=0.5mm](r) edge node[above] {$\epsilon$} (q2);
				\path[->,line width=0.5mm](r) edge [loop above] node[left] {$\gamma$} (r);
			\end{tikzpicture}
		\end{itemize}
	\end{frame}
	
	\begin{frame} \frametitle{Aufgabe 3}
		\small
		Geben Sie zu jedem der regulären Ausdrücke $r_i$ einen 
		NFA $\mathcal M_i$ mit $L(\mathcal M_i) = L(r_i)$ an.
		
		\begin{enumerate}[a)]
			\item $r_1 = (ab)^\ast$
			\item $r_2 = a(b \mid c)a^\ast \mid a^\ast$
			% \item $r_3 = (bb\mid cc^\ast)^\ast$
		\end{enumerate}
		
		Wenden Sie dabei jeweils den \textit{kompositionellen Ansatz} sowie den \textit{expliziten Ansatz} zur Konstruktion von NFAs aus der Vorlesung an.
	\end{frame}

	\section{Aufgabe 4: \\ \itshape NFA $\to$ RegExp}


	\begin{frame}\frametitle{NFA $\to$ RegExp: Ersetzungsmethode}
		\small
		\textbf{Gegeben}: NFA  $\Smach{M}=\tuple{Q,\Sigma,\delta,Q_0,F}$
		
		\textbf{Gesucht}: regulärer Ausdruck $\alpha$ mit $\Slang{L}(\alpha)=\Slang{L}(\Smach{M})$
		
		\textbf{Idee}:
		
		\codebox{Für jeden Zustand $q\in Q$, berechne einen regulären Ausdruck $\alpha_q$ für die Sprache $\Slang{L}(\alpha_q) = \Slang{L}(\Smach{M}_q)$ mit $\Smach{M}_q=\tuple{Q,\Sigma,\delta,\{q\},F}$}
		
		Für $Q_0 = \menge{q_1,q_2,\dots, q_n}$ gilt dann
		\begin{equation*}
			\Slang{L}(\Smach{M})
			=\bigcup_{q\in Q_0} \Slang{L}(\alpha_q)
			=  \Slang{L}(\alpha_{q_1}\mid\alpha_{q_2}\mid\dots\mid \alpha_{q_n})
		\end{equation*}
	\end{frame}

	\begin{frame}		
		\small
		\begin{enumerate}[(1)]
			\item \textbf{Vereinfache den Automaten} (entferne offensichtlich unnötige Zustände)	
			
			\item \textbf{Bestimme das Gleichungssystem} 
			
			\textit{Intuition}: Beschreibe $\alpha_q$ in Abhängigkeit von Folgezuständen
			
			\defbox{ \\[-0.9\baselineskip]
				\begin{itemize}
					\item Für jeden Zustand $q\in Q \setminus F$:
					$\alpha_q \equiv \sum_{\Sterm{a}\in\Sigma}\; \sum_{p\in\delta(q,\Sterm{a})} \Sterm{a}\alpha_p$
					\item Für jeden Zustand $q\in F$:
					$\alpha_q \equiv \bluealert{\epsilon} \mid \sum_{\Sterm{a}\in\Sigma}\; \sum_{p\in\delta(q,\Sterm{a})} \Sterm{a}\alpha_p$
				\end{itemize}
			}
			%
			\item \textbf{Löse das Gleichungssystem} durch Einsetzen und
			
			\theobox{Regel von Arden: Aus $\alpha\equiv \beta\alpha\mid\gamma$ mit $\epsilon\notin\Slang{L}(\beta)$ folgt $\alpha\equiv\beta^\ast\gamma$.}
			%
			\item \textbf{Gib den Ausdruck für die Sprache des NFA an}
			
			Für $Q_0 = \menge{q_1,q_2,\dots, q_n}$ gilt dann
			\begin{equation*}
				\Slang{L}(\Smach{M})
				=\bigcup_{q\in Q_0} \Slang{L}(\alpha_q)
				=  \Slang{L}(\alpha_{q_1}\mid\alpha_{q_2}\mid\dots\mid \alpha_{q_n})
			\end{equation*}
		\end{enumerate}
	\end{frame}

	\newcommand{\Lijk}[3]{\ensuremath{\Slang{L}^{#3}[#1,#2]}}
	\newcommand{\alphaijk}[3]{\ensuremath{\alpha^{#3}[#1,#2]}}
	
	\begin{frame} \frametitle{NFA $\to$ RegExp: Dynamische Ermittlung}
		\small
		\textbf{Gegeben}: NFA  $\Smach{M}=\tuple{Q,\Sigma,\delta,Q_0,F}$
		
		\textbf{Gesucht}: regulärer Ausdruck $\alpha$ mit $\Slang{L}(\alpha)=\Slang{L}(\Smach{M})$
		
		\textbf{Ansatz}:
		
		\codebox{Für jedes Paar von Zuständen $q,p\in Q$, berechne einen regulären Ausdruck $\alpha_{q,p}$ für die Sprache
		$\Slang{L}(\alpha_{q,p}) = \Slang{L}(\Smach{M}_{q,p})$ mit $\Smach{M}_{q,p}=\tuple{Q,\Sigma,\delta,\{q\},\{p\}}$}
		
		Dann gilt:
		\begin{align*}
			\Slang{L}(\Smach{M}) &=\bigcup_{q\in Q_0} \bigcup_{p\in F} \Slang{L}(\alpha_{q,p})
			=  \Slang{L}\left(\sum_{q\in Q_0} \sum_{p\in F} \alpha_{q,p}\right)
		\end{align*}
	
		\begin{itemize}
			\item $\Lijk{i}{j}{k} \dots$ Sprache mit Start in $q_i$, Ende in $q_j$ und nutzt nur Zwischenzustände $q_1, \dots, q_k$
			\item $\alphaijk{i}{j}{k}$ zugehöriger regulärer Ausdruck
		\end{itemize}
	\end{frame}
	
	\begin{frame}
		\small
		\textbf{Idee}: lasse immer mehr Zwischenzustände von $q$ nach $p$ zu \\
		\begin{itemize}
			\item $k = n$: nutze alle Zustände --- Ergebnis ablesbar
			\item $k = 0$: nutze keine Zwischenzustände --- $\alpha^0[i,j]$ direkt ablesbar:
			\codebox{
				Sei $\{\Sterm{a_1},\ldots,\Sterm{a_m}\}=\{\Sterm{a}\in\Sigma\mid q_i\stackrel{\Sterm{a}}{\to}q_j\}$ die Menge der Beschriftungen von direkten Übergängen von $q_i$ zu $q_j$.
				\begin{itemize}
					\item Falls $i\neq j$, dann ist
					$\alphaijk{i}{j}{0} = \Sterm{a_1}\mid\ldots\mid\Sterm{a_m}$
					\item Falls $i=j$, dann ist
					$\alphaijk{i}{j}{0} = \Sterm{a_1}\mid\ldots\mid\Sterm{a_m}\mid\epsilon$
				\end{itemize}
			}
		\end{itemize}
	
		\textbf{Update-Formel}:
		\codebox{\vspace{-1ex}
			\begin{equation*}
				\alphaijk{i}{j}{k+1} = \alphaijk{i}{j}{k} \mid \big( \textcolor{darkred}{\alphaijk{i}{k+1}{k}} (\textcolor{darkblue}{\alphaijk{k+1}{k+1}{k}})^\ast \textcolor{darkgreen}{\alphaijk{k+1}{j}{k}} \big)
			\end{equation*}
		}
	
		\begin{itshape}
			\scriptsize
			vgl. VL \enquote{Algorithmen \& Datenstrukturen}, Prozess-Problem im Aho-Hopcroft-Ullman-Algorithmus
		\end{itshape}
	\end{frame}

	\begin{frame} \frametitle{Aufgabe 4}
		\small
		Entwickeln Sie für die Sprache $L$ über dem Alphabet $\Sigma = \menge{a, b, c}$ einen regulären Ausdruck $r$ mit $L = L(r)$. Für alle Wörter $w \in L$ gilt:
		\begin{itemize}
			\item $w$ enthält $aaa$.
			\item $w$ endet mit $c$.
			\item Die Anzahl der $b$ in $w$ ist gerade.
		\end{itemize}
	\end{frame}
	
	\begin{frame}
		\small
		\textbf{Idee}:
		$L = 
		\underbrace{\brackets{\menge{w \in \Sigma^\ast : w \text{ enthält } aaa} \cap \menge{w \in \Sigma^\ast : \abs{w}_b \text{ gerade} }}}_{\defeq L(\mathcal{M})} \cdot \menge{c}$
	
		\begin{itemize}
			\item Konstruiere Automaten für die erste und dritte Bedingung
			\item Produktautomat $\mathcal{M}$ erkennt die Sprache ohne $c$ am Ende
		\end{itemize}
	
		\begin{center}
			\scalebox{0.6}{
			\begin{tikzpicture}[%
				->,
				>=stealth',
				semithick,
				initial text=,
				shorten <=2pt,   
				shorten >=2pt,
				auto, 
				on grid,
				node distance=14ex and 7em,
				every state/.style={minimum size=0pt,inner sep=2pt,text height=1.5ex,text depth=.25ex},
				bend angle=20]
				\node                  (init)                 {$\mathcal{M}:$};
				\node[state,initial]   (q0p0) [right=of init] {$\tuple{q_0, p_0}$};
				\node[state]           (q1p0) [right=of q0p0] {$\tuple{q_1, p_0}$};
				\node[state]           (q2p0) [right=of q1p0] {$\tuple{q_2, p_0}$};
				\node[state,accepting] (q3p0) [right=of q2p0] {$\tuple{q_3, p_0}$};
				\node[state]           (q0p1) [below=of q0p0] {$\tuple{q_0, p_1}$};
				\node[state]           (q1p1) [right=of q0p1] {$\tuple{q_1, p_1}$};
				\node[state]           (q2p1) [right=of q1p1] {$\tuple{q_2, p_1}$};
				\node[state]           (q3p1) [right=of q2p1] {$\tuple{q_3, p_1}$};
				
				\path[->,draw] {
					(q0p0) edge[loop above] node {$a,c$} (q0p0)
					       edge node {$a$}   (q1p0)
					       edge[bend left] node {$b$}   (q0p1)
					(q1p0) edge node {$a$}   (q2p0)
					(q2p0) edge node {$a$}   (q3p0)
					(q3p0) edge[loop above] node {$a,c$} (q3p0)
					(q3p0) edge[bend left] node {$b$}   (q3p1)
					(q0p1) edge[loop below] node {$a,c$} (q0p1)
					       edge[bend left] node {$b$} (q0p0)
					(q0p1) edge node {$a$}   (q1p1)
					(q1p1) edge node {$a$}   (q2p1)
					(q2p1) edge node {$a$}   (q3p1)
					(q3p1) edge[loop below] node {$a,c$} (q3p1)
					       edge[bend left] node {$b$}   (q3p0)
				};
			\end{tikzpicture}}
		\end{center}
		\vspace{-\baselineskip}
	\end{frame}

	\begin{frame}
		\small
		\begin{itemize}
			\item Umwandlung in regulären Ausdruck $\alpha_{\mathcal{M}}$ (z.B. via Lemma von Arden)		
		\end{itemize}
		\begin{align*}
			\alpha_{\mathcal{M}} = \
			&\Big( a \mid c \mid  b \, (a | c)^\ast \, b \Big)^\ast \\
			&\Big( aaa ( a \mid c \mid  b \, (a | c)^\ast \, b )^\ast
			\bigmid
			b (a|c)^\ast aaa(a|c)^\ast b \, \big(a \mid c \mid b(a|c)^\ast b \big)^\ast
			\Big)
		\end{align*}
		\begin{itemize}
			\item Mit $r = \alpha_{\mathcal{M}} c$ gilt dann $L(r) = L$.
		\end{itemize}		
	\end{frame}
	
	\section{Aufgabe 5: \\ \itshape Minimierung von Automaten}

	\begin{frame}\frametitle{Äquivalenz von Zuständen \& Quotientenautomat}
		\small
		DFA $\Smach{M}=\tuple{Q,\Sigma,\delta,q_0,F} \leadsto$ DFA $\Smach{M}_q=\tuple{Q,\Sigma,\delta,q,F}$
		
		\defbox{
			\makebox[5.5cm][l]{\textbf{Äquivalenz von Zuständen}:} $p\sim_{\Smach{M}}q 
			\quad \Leftrightarrow \quad
			\Slang{L}(\Smach{M}_{p}) = \Slang{L}(\Smach{M}_{q})$ \\
			\makebox[5.5cm][l]{\textbf{Äquivalenzklasse}:} $[q]_{\sim}=\{p\in Q\mid q\sim p\}$ \\
			\makebox[5.5cm][l]{\textbf{Quotient} von $P \subseteq Q$:} $\simquot{P\!} = \menge{[p]_{\sim}\mid p \in P}$
		}
	
	
		\textbf{Quotientenautomat}:  Verschmelzen von äquivalenten Zuständen
		
		\defbox{Für einen DFA $\Smach{M} = \tuple{Q,\Sigma,\delta,q_0,F}$ mit totaler Übergangsfunktion ist der Quotientenautomat $\simquot{\Smach{M}} =\tuple{\simquot{Q},\Sigma,\delta_{\sim},[q_0]_{\sim_{\Smach{M}}},\simquot{F\!}}$ gegeben durch
		\begin{itemize}
			\item $\simquot{Q} = \{[q]_{\sim}\mid q\in Q\}$
			\item $\delta_{\sim}([q]_{\sim},\Sterm{a}) = [\delta(q,\Sterm{a})]_{\sim}$
			\item $\simquot{F\!} = \{[q]_{\sim}\mid q\in F\}$
		\end{itemize}
		}
	\end{frame}

	\begin{frame}
		\small
		
		\textbf{Bestimmung von $\sim$}:
		\vspace{-.7\baselineskip}
		\begin{itemize}
			\item Initialisiere ${\not\sim}\defeq\emptyset$
			\item \textbf{Regel 1}: Für jedes Paar von Zuständen $\tuple{q,p}\in Q\times Q$:\\
			falls $q\in F$ und $p\notin F$, dann "`speichere $q\not\sim p$"'
			% ${\not\sim}\defeq{\not\sim}\cup\{\tuple{q,p},\tuple{p,q}\}$
			\item \textbf{Regel 2}: Für jedes Paar $\tuple{q,p}\in Q\times Q\setminus {\not\sim}$ und jedes $\Sterm{a}\in\Sigma$:\\
			falls $\delta(q,\Sterm{a})\not\sim\delta(p,\Sterm{a})$ dann "`speichere $q\not\sim p$"'
			% 	${\not\sim}\defeq{\not\sim}\cup\{\tuple{q,p},\tuple{p,q}\}$
			\item Wiederhole Regel 2 bis keine Änderungen mehr auftreten
			\item Das Ergebnis ist $(Q\times Q)\setminus {\not\sim}$
		\end{itemize}
		
		\examplebox{
			\begin{minipage}{6cm}
				\begin{flushleft}
					\textbf{Beispiel}: Für einen DFA mit Zuständen $Q=\{A,B,C,D,E\}$ genügt eine Tabelle mit zehn Feldern (statt $5^2=25$).\\[1ex]
					
					{\tiny (dazu reihen wir Zustände vertikal in umgekehrter Reihenfolge)}
				\end{flushleft}
			\end{minipage}
			\narrowcentering{\scalebox{0.8}{%
					\begin{tabular}{c|c|c|c|c|}
						\multicolumn{1}{c}{}  & \multicolumn{1}{c}{A} & \multicolumn{1}{c}{B} & \multicolumn{1}{c}{C} & \multicolumn{1}{c}{D}  \\\cline{2-5}
						E &   &   &   &    \\\cline{2-5}
						D &   &   &        \\\cline{2-4} 
						C &   &            \\\cline{2-3}
						B &                \\\cline{2-2}
			\end{tabular}}}
		}
	\end{frame}
	
	\begin{frame} \frametitle{Aufgabe 5}
		\small
		Berechnen Sie für folgenden DFA $\mathcal{M}=(\{q_0,q_1,q_2,q_3,q_4,q_5\},\{a,b\},\delta,q_0,\{q_1,q_2,q_4\})$ mit $\delta$:
		
		\begin{center}
			\begin{tikzpicture}[%
				->,
				>=stealth',
				semithick,
				initial text=,
				shorten <=2pt,   
				shorten >=2pt,
				auto, 
				on grid,
				node distance=14ex and 8em,
				every state/.style={minimum size=0pt,inner sep=2pt,text height=1.5ex,text depth=.25ex},
				bend angle=20]
				\path {
					node[state,initial] (q_0) {$q_0$}
					node[state,accepting] (q_1) [right=of q_0] {$q_1$}
					node[state,accepting] (q_2) [right=of q_1] {$q_2$}
					node[state] (q_3) [below=of q_0] {$q_3$}
					node[state,accepting] (q_4) [right=of q_3] {$q_4$}
					node[state] (q_5) [right=of q_4] {$q_5$}
				};
				
				\path[->,draw] {
					(q_0) edge node {$b$} (q_1)
					(q_1) edge node {$a$} (q_2)
					(q_0) edge[bend left] node {$a$} (q_3)
					(q_3) edge[bend left] node {$a$} (q_0)
					(q_3) edge node {$b$} (q_4)
					(q_4) edge node {$b$} (q_5)
					(q_4) edge node[pos=0.3] {$a$} (q_2)
					(q_1) edge node[pos=0.7] {$b$} (q_5)
					(q_2) edge node {$b$} (q_5)
					(q_2) edge[loop right] node {$a$} (q_2)
					(q_5) edge[loop right] node {$a,b$} (q_5)
				};
				
			\end{tikzpicture}
		\end{center}
		die Äquivalenzrelation
		$\sim_{\mathcal{M}}$, und geben Sie den Quotientenautomaten ${\mathcal{M}}/_\sim$ an.
	\end{frame}

	\begin{frame}
		\small
		\begin{alignat*}{2}
			\sim &=& &\left\{
				\tuple{q_1, q_2}, \tuple{q_2, q_4}, \tuple{q_1, q_4}, \tuple{q_0, q_3}  \right\} \\
				&\cup & &\left\{ \tuple{q_2, q_1}, \tuple{q_4, q_2}, \tuple{q_4, q_1}, \tuple{q_3, q_0} \right\} \\
				&\cup& &\menge{\tuple{q,q} : q \in Q} \\
			Q/_\sim &=& &\Big\{
			\underbrace{\menge{q_0, q_3}}_{ = [q_0]_\sim},
			\underbrace{\menge{q_1, q_2, q_4}}_{ = [q_1]_\sim},
			\underbrace{\menge{q_5}}_{= [q_5]_\sim}%
			\Big\}
		\end{alignat*}
	
		\begin{center}
			\scalebox{1}{
				\begin{tikzpicture}[%
					->,
					>=stealth',
					semithick,
					initial text=,
					shorten <=2pt,   
					shorten >=2pt,
					auto, 
					on grid,
					node distance=14ex and 8em,
					every state/.style={minimum size=0pt,inner sep=2pt,text height=1.5ex,text depth=.25ex},
					bend angle=20]
					\node[state,initial]   (q03)                  {$\menge{q_0, q_3}$};
					\node[state,accepting] (q124) [right=of q03]  {$\menge{q_1, q_2, q_4}$};
					\node[state]           (q5)   [right=of q124] {$\menge{q_5}$};
					
					\path[->,draw] {
						(q03)  edge[loop above] node {$a$}   (q03)
						       edge             node {$b$}   (q124)
						(q124) edge[loop above] node {$a$}   (q124)
						       edge             node {$b$}   (q5)
						(q5)   edge[loop above] node {$a,b$} (q5)
					};
			\end{tikzpicture}}
		\end{center}
	\end{frame}
	

	
\end{document}